\section{Occidental Petroleum (NYSE: OXY)}
\label{sec:oxy}

Kurs \je{\OxyKurs}{USD} am \OxyKursDatum{}\QW{Yahoo Finance, Kursreihe OXY}{quelle:kurse}{19.09.2026},
\Mio{\OxyAktienMio}{} Aktien\Q{OxyZK}{Deckblatt}, B\"orsenwert \Mio{\OxyBoersenwert}{USD}.

\subsection{Eigent\"umer und Directors' Dealings}

Warren E.~Buffett und verbundene Gesellschaften (Berkshire Hathaway) halten
\Pz{\OxyBerkshireAnteil} der Stammaktien, einschlie{\ss}lich
\Mio{\OxyBerkshireOptionsscheine}{} Aktien aus Optionsscheinen; es folgen
Dodge~\&~Cox mit \Pz{\OxyDodgeAnteil} und The Vanguard Group mit
\Pz{\OxyVanguardAnteil}\Q{OxyPS}{65}. Die Bilanz f\"uhrt Vorzugsaktien von
\Mio{\OxyVorzugsaktien}{USD} im Eigenkapital von \Mio{\OxyEigenkapital}{USD}\Q{OxyZK}{59}.
Berkshire hat zudem das Chemiegesch\"aft OxyChem f\"ur rund \je{\OxyChemPreis}{Mrd.~USD}
erworben; der Vollzug war am 2.~Januar 2026\Q{OxyPS}{32}.

Directors' Dealings der letzten zw\"olf Monate (SEC-Formular~4): \OxyDdKaeufe{} K\"aufe am
offenen Markt, \OxyDdVerkaeufe{} Verk\"aufe. William R.~Klesse (Director) kaufte
\OxyKlesseStueck{} Aktien zu \je{\OxyKlessePreis}{USD}, Richard A.~Jackson (President und
CEO) \OxyJacksonStueck{} Aktien zu \je{\OxyJacksonPreis}{USD}\QW{SEC EDGAR, Formulare 4 der
Occidental Petroleum}{quelle:form4}{19.09.2026}. Zuteilungen, Steuereinbehalte und
Schenkungen sind Verg\"utungsvorg\"ange und nicht mitgez\"ahlt.

\bk{Ein Drittel der Stimmen in der Hand eines Aktion\"ars, der zugleich Optionsscheine
h\"alt und gerade ein Gesch\"aftsfeld vom Unternehmen gekauft hat, ist mehr als eine
Beteiligung: Er ist Aktion\"ar und Vertragspartner zugleich. Wer die Vorzugsaktien h\"alt,
nennen die hier gelesenen Seiten nicht. Zwei K\"aufe der Organe sind ein Datenpunkt, kein
Signal.}

\subsection{Wettbewerbsposition}

\begin{table}[htbp]\centering\small
\caption{Behauptete Vorteile und die Zahl, die sie zeigen w\"urde (Occidental).}\label{tab:moat-oxy}
\begin{tabularx}{\linewidth}{@{}XXl@{}}\toprule
Behauptung (Quelle) & Zahl, die sie zeigen w\"urde & H\"alt?\\\midrule
\enquote{ranks among the largest oil and gas producers in the U.S.}\Q{OxyZK}{3} &
F\"orderanteil an der US-Produktion & nicht belegbar\\
\enquote{leading positions in the Permian and DJ Basins as well as offshore Gulf of
America}\Q{OxyZK}{3} & Anteil an der F\"orderung je Becken & nicht belegbar\\
\enquote{the largest independent oil producer in Oman}\Q{OxyZK}{3} & F\"orderung gegen
andere unabh\"angige Produzenten & nicht belegbar\\
\enquote{industry leader in Permian Basin CO$_2$ flooding}\Q{OxyZK}{26} & Zahl der Flutungen
gegen Wettbewerber & teilweise\\
\bottomrule\end{tabularx}\end{table}

Der 10-K nennt die Zahl der eigenen CO$_2$-Flutungen, aber keine Vergleichszahl eines
Wettbewerbers; die \"ubrigen Behauptungen stehen ohne Zahl.

\bk{Ein F\"orderer von Rohstoffen verkauft zum Weltmarktpreis. Seine Position zeigt sich in
den F\"orderkosten je Barrel, nicht in einer Rangliste, und die Kosten je Barrel im Vergleich
zu Wettbewerbern liegen in den hier gelesenen Quellen nicht vor.}

\subsection{Mehrjahresreihe}

\reihentabelle{oxy}{USD}{oxy}

Der freie Mittelzufluss lag \OxyFcfJahre{} im Median bei \Mio{\OxyFcfMedian}{USD}; 2025
betrug er \Mio{\OxyFcfVoll}{USD} und lag damit \"uber \Pz{\OxyFcfRang} der Jahre der eigenen
Reihe. Ab 2023 ist der freie Zufluss ohne OxyChem gerechnet, die Spalte
\enquote{Operativ} zeigt weiter den gesamten Konzern: Der 10-K weist f\"ur 2025 einen
operativen Mittelzufluss der fortgef\"uhrten T\"atigkeit von \Mio{\OxyCfoFortgefuehrt}{USD}
und der aufgegebenen von \Mio{\OxyCfoAufgegeben}{USD} aus\Q{OxyZK}{63}.

\annahme{Der 10-K trennt die Investitionen nicht nach T\"atigkeit. F\"ur 2023 bis 2025 sind
die Investitionen der fortgef\"uhrten T\"atigkeit als Investitionen gesamt abz\"uglich des
Investitions-Mittelflusses der aufgegebenen T\"atigkeit gerechnet. Vor 2023 enth\"alt die
Reihe OxyChem; sie ist dort mit dem heutigen Unternehmen nicht deckungsgleich.}

\subsection{Kursverlauf}

\kgvtabelle{oxy}{USD}

\bk{Der Kurs folgt nicht dem Ergebnis: Das KGV liegt heute bei \OxyKgvHeute{} gegen einen
Median von \OxyKgvMedian{} (\OxyKgvJahre, nur Jahre mit Gewinn) und damit h\"oher als in
\Pz{\OxyKgvRang} der erfassten Jahre. Gesunken ist das Ergebnis je Aktie, nicht der Kurs;
der Markt bewertet einen \"Olpreiszyklus, nicht das Ergebnis eines Jahres.}

\subsection{Votum}

\schwellentabelle{oxy}{USD}{oxy}
\umkehrtabellen{oxy}

\vd{\textbf{Nicht kaufen} zu \je{\OxyKurs}{USD}. Die Schwelle liegt auf zehn Jahre bei
\je{\OxySchwelleVOLLZehn}{USD}, wenn das Niveau von 2025 gehalten wird, und bei
\je{\OxySchwelleMEDZehn}{USD} auf dem Median der eigenen Reihe; der Kurs ist das
\OxyVielfachesVOLL-fache der ersten. F\"ur den Halter: kein Nachkauf; ein Verkauf
w\"are eine Wette auf den \"Olpreis, die diese Rechnung weder st\"utzt noch widerlegt.
F\"ur den Nichthalter: warten.}

\vd{\emph{Begr\"undung.} Der heutige Preis setzt ein Wachstum des freien Zuflusses von
\OxyWachstumVOLLZehn{} p.\,a.\ \"uber zehn Jahre voraus, oder einen Endwert vom
\OxyEndwertVOLLZehn-fachen des Buchwerts. Die realisierte Rate von \OxyFcfWachstumReal{}
(\OxyFcfWachstumSpanne) ist kein Gegenbeweis: Sie misst die Erholung aus dem
\"Olpreistief 2015/16 und nicht ein Wachstum des Gesch\"afts. Auch zur niedrigeren H\"urde
von \Pz{\HuerdeLang} bleibt die Schwelle mit \je{\OxySchwelleVOLLZehnLang}{USD} unter dem
Kurs.}

\vd{\emph{Was das Votum \"andert.} Ein Kurs unter der Schwelle; ein Jahr, in dem der freie
Zufluss nach Vorzugsdividende \"uber dem von 2025 liegt, ohne dass der \"Olpreis steigt; eine
Tilgung der Vorzugsaktien, die den Zufluss der Stammaktion\"are erh\"oht.}

\vd{\emph{Was dagegen spricht.} Das Modell h\"alt den \"Olpreis auf dem Niveau des
jeweiligen Szenarios fest und gibt dem Anteil am Ende nur den Buchwert; steigt der
\"Olpreis, steigt der Zufluss mit, und der Endwert eines F\"orderers mit Reserven liegt
selten beim Buchwert. Wer eine eigene Einsch\"atzung zum \"Olpreis hat, rechnet mit
einem anderen Zufluss; dieses Dokument hat keine.}
